\title{DeepSeekMath: Pushing the Limits of Mathematical Reasoning in Open Language Models}

\begin{abstract}
Mathematical reasoning remains a formidable hurdle for language models, stemming from its inherent complexity and rigid structure. Here, we present DeepSeekMath~7B, developed by further pre-training DeepSeek-Coder-Base-v1.5 7B on a corpus of 120B tokens rich in mathematical content, extracted from Common Crawl and supplemented with data in natural language and code. Without resorting to external toolchains or ensemble-based voting, DeepSeekMath~7B attains a notable 51.7\% on the competition-grade MATH benchmark, coming close to the results achieved by Gemini-Ultra and GPT-4. Moreover, using self-consistency over 64 generated samples, DeepSeekMath~7B reaches 60.9\% on MATH. The substantial progress in the mathematical reasoning abilities of DeepSeekMath~7B can be traced to two primary innovations: first, we tap into the immense utility of openly available web resources via a carefully crafted data selection procedure; second, we propose Group Relative Policy Optimization (GRPO)—an adaptation of Proximal Policy Optimization (PPO)—to both improve mathematical reasoning and reduce the memory footprint of PPO.
\end{abstract}

\section{Introduction}

Large language models (LLMs) have fundamentally transformed artificial intelligence approaches to mathematical reasoning, yielding major breakthroughs in quantitative reasoning assessments \citep{MATH} as well as in geometric problem-solving tasks \citep{trinh2024solving}. These systems are also increasingly capable of providing meaningful assistance to humans tackling difficult mathematical challenges \citep{tao}. Nevertheless, the latest high-performing systems, including GPT-4 \citep{gpt4} and Gemini-Ultra \citep{gemini}, remain inaccessible to the public, while the open-source alternatives still significantly lag in terms of mathematical reasoning proficiency.

This work presents DeepSeekMath, a language model tailored to the mathematical domain that not only surpasses the open-source alternatives in mathematical reasoning but also comes close to matching GPT-4 on standardized academic assessments. Central to our approach is the construction of the DeepSeekMath Corpus, a vast and meticulously curated pre-training dataset consisting of 120B mathematics-focused tokens. The dataset is sourced from Common Crawl (CC) by deploying a classifier grounded in fastText \citep{joulin2016fasttext}. During the first stage, this classifier is trained with examples from OpenWebMath \citep{openwebmath} as the positive class, and various other internet pages as negative samples. Using this model, we proceed to extract further positive cases from CC, which are subsequently validated and refined via human annotation. The improved dataset is then utilized to retrain the classifier for heightened accuracy. The performance of our base model, DeepSeekMath-Base 7B, illustrates the corpus’s quality, recording scores of 64.2\% on GSM8K \citep{gsm8k} and 36.2\% on the advanced MATH benchmark \citep{MATH}, outperforming even Minerva 540B \citep{minerva}. Notably, since the DeepSeekMath Corpus encompasses multiple languages, we observe marked gains on Chinese mathematical evaluations \citep{wei2023cmath,agieval}. We suggest that our methodology for mathematical data processing can provide a solid foundation for subsequent work, leaving ample opportunities for future advancement in this area.

DeepSeekMath-Base is constructed by initializing with DeepSeek-Coder-Base-v1.5 7B \citep{deepseek-coder}, given our observation that leveraging a code-focused pretraining model yields superior results over starting from a standard large language model (LLM). Additionally, we find that mathematical training not only strengthens the model's mathematical proficiency but also enhances its general reasoning skills, as evidenced by improvements on MMLU \citep{mmlu} and BBH benchmarks \citep{bbh}.

Upon completing pre-training, we perform mathematical instruction tuning on DeepSeekMath-Base, incorporating datasets designed for chain-of-thought \citep{cot}, program-of-thought \citep{pot,pal}, and tool-integrated reasoning \citep{tora}. The resulting DeepSeekMath-Instruct 7B model outperforms all existing 7B parameter models and achieves results on par with larger 70B open-source instruction-tuned models.

In addition, we present Group Relative Policy Optimization (GRPO), an alternative reinforcement learning (RL) algorithm inspired by Proximal Policy Optimization (PPO) \citep{schulman2017proximal}. Distinctively, GRPO eliminates the need for a critic network by estimating baselines through group scoring, leading to significant reductions in computational costs. With the use of only a portion of English instruction tuning data, GRPO produces marked gains over the already strong DeepSeekMath-Instruct model in both domain-specific (GSM8K: 82.9\% $\rightarrow$ 88.2\%, MATH: 46.8\% $\rightarrow$ 51.7\%) and out-of-domain mathematical benchmarks (such as CMATH: 84.6\% $\rightarrow$ 88.8\%) during RL training. 

Furthermore, we establish a comprehensive framework to interpret a range of methods, including Rejection Sampling Fine-Tuning (RFT) \citep{yuan2023scaling}, Direct Preference Optimization (DPO) \citep{dpo}, PPO, and GRPO, all of which can be viewed as forms of either direct or simplified reinforcement learning within this unified schema. We also conduct systematic experiments—examining dimensions such as online versus offline training, process versus outcome supervision, and single-turn versus iterative RL—to thoroughly analyze the critical factors underpinning this framework. Finally, we offer insights into the mechanisms by which RL enhances the effectiveness of instruction-tuned models and outline prospective avenues to further advance RL methods using this integrative approach.

\subsection{Contributions}
This work presents advancements in scalable mathematical pre-training, complemented by investigations into reinforcement learning methodologies.

\textbf{Math Pre-Training at Scale}

\begin{itemize}[topsep=0pt]
    \item We demonstrate that publicly available Common Crawl data holds substantial utility for mathematical tasks. Through a carefully engineered data selection framework, we assemble the DeepSeekMath Corpus—a rigorously curated dataset comprising 120B tokens from web pages rich in mathematical content. This collection is nearly 7 times greater than the math web page dataset leveraged by Minerva \citep{minerva}, and about 9 times larger than the more recent OpenWebMath \citep{openwebmath}.

    \item Our DeepSeekMath-Base 7B model, trained on this dataset, matches the performance level of Minerva 540B \citep{minerva}, suggesting that the total parameter count is not the predominant factor in mathematical reasoning strength. High-caliber data enables even smaller models to attain impressive capabilities.

    \item Experimental results from our math training regimen indicate that pre-training on code prior to math-specific data enhances the models’ proficiency at solving mathematical questions, both with and without external tool assistance. This partially addresses the ongoing debate on whether code training boosts reasoning skill: our findings affirm its positive impact, particularly for mathematical inference.

    \item While it is standard practice to train on arXiv papers, we observe that such pre-training does not yield significant performance gains across any of the mathematical benchmarks evaluated in this study.
\end{itemize}

\textbf{Exploration and Analysis of Reinforcement Learning}

\begin{itemize}[topsep=0pt]
    \item We present Group Relative Policy Optimization (GRPO), a reinforcement learning approach designed to be both resource-efficient and performant. Unlike Proximal Policy Optimization (PPO), GRPO eliminates reliance on a critic model by leveraging group scores to estimate the baseline, thereby substantially lowering computational demands during training.
    \item Through empirical evaluation, we show that applying GRPO to our instruction-tuned system, DeepSeekMath-Instruct, leads to significant gains when using only instruction-tuning data. We also note marked improvements in handling out-of-domain tasks during the reinforcement learning phase.
    \item We establish a cohesive analytical framework that contextualizes various approaches, including RFT, DPO, PPO, and GRPO. Additionally, we undertake comprehensive experimental investigations—such as comparing online versus offline learning strategies, outcome versus process-oriented supervision, and single-turn versus iterative reinforcement learning—to thoroughly examine critical aspects of this unified perspective.
    \item Building upon this shared framework, we analyze why reinforcement learning proves effective in this domain and highlight several avenues for further advancements in optimizing large language model reinforcement learning.
\end{itemize}

\subsection{Summary of Evaluations and Metrics}

\begin{itemize}[topsep=0pt]
    \item \textbf{English and Chinese Mathematical Reasoning}: We systematically evaluate our models on a range of mathematical reasoning benchmarks in both English and Chinese, encompassing tasks from primary through collegiate difficulty. English benchmarks utilized include GSM8K \citep{gsm8k}, MATH \citep{MATH}, SAT \citep{llemma}, OCW Courses \citep{minerva}, and MMLU-STEM \citep{mmlu}; for Chinese, we rely on MGSM-zh \citep{mgsm}, CMATH \citep{wei2023cmath}, Gaokao-MathCloze \citep{agieval}, and Gaokao-MathQA \citep{agieval}. Our assessments cover models’ capacities to produce comprehensive, tool-free text solutions, as well as solve tasks programmatically via Python.

    DeepSeekMath-Base demonstrates comparable performance to the proprietary Minerva 540B \citep{minerva} model on English benchmarks, and it clearly outperforms existing open-source base models, including Mistral 7B \citep{mistral} and Llemma-34B \citep{llemma}, irrespective of their mathematical pre-training history—often by wide margins. Particularly, DeepSeekMath-Base leads on Chinese evaluations, a result likely attributed to our approach of supplementing English data with high-quality non-English mathematical corpora, as opposed to prior studies \citep{minerva,llemma} that restrict data collection to English. With the addition of instruction-based tuning and reinforcement learning, DeepSeekMath-Instruct and DeepSeekMath-RL reach notable performance, crossing the 50\% accuracy threshold on the rigorous, competition-level MATH benchmark—an open-source first.
\end{itemize}

\item \textbf{Formal Mathematics}: DeepSeekMath-Base is assessed on the informal-to-formal theorem proving task presented in \citep{dsp_proof}, using miniF2F \citep{minif2f} as the testbed and Isabelle \citep{isabelle} as the proof assistant. The results indicate that DeepSeekMath-Base achieves notable performance in few-shot autoformalization settings.
\item \textbf{Natural Language Understanding, Reasoning, and Code}: To provide a thorough evaluation of the models’ abilities in comprehension, reasoning, and programming, we test DeepSeekMath-Base on the Massive Multitask Language Understanding (MMLU) benchmark \citep{mmlu}, which includes 57 multiple-choice questions spanning a broad array of topics. Additionally, we use BIG-Bench Hard (BBH) \citep{bbh}, featuring 23 multi-step reasoning tasks, alongside HumanEval \citep{codex} and MBPP \citep{mbpp}—both commonly used to benchmark code generation. Our results show that pre-training on mathematical data enhances both reasoning and general language understanding.

\section{Math Pre-Training}

\subsection{Data Collection and Decontamination} This section details the assembly of the DeepSeekMath Corpus from Common Crawl. As illustrated in Figure \ref{fig:data_collect}, our iterative pipeline demonstrates a systematic approach for extracting a large-scale mathematics-focused dataset from Common Crawl, starting from an initial high-quality, math-centric seed dataset. While the following methodology is demonstrated for mathematics, it can readily be applied to other specialized domains such as computer code.

Initially, we select OpenWebMath \citep{openwebmath}—a curated collection of mathematical web texts—as the seed corpus. Utilizing this resource, we train a fastText model \citep{joulin2016fasttext} to identify additional web pages similar to those in OpenWebMath. Concretely, 500,000 examples are sampled at random from the seed as positive instances, while an equal number of pages from Common Crawl serve as negatives. The model is trained with the following parameters using an open-source implementation\footnote{\url{https://fasttext.cc}}: a vector size of 256, learning rate of 0.1, n-grams with a maximum length of 3, minimum occurrence count per word set to 3, and 3 epochs. To decrease the scale of the original Common Crawl dataset, we implement deduplication both at the URL level and through near-duplicate detection, ultimately reducing it to 40B HTML web pages. These are then scored for mathematical content with our fastText model, and only the highest-ranking pages—judged by the model’s predictions—are retained, thereby filtering out irrelevant or low-quality content. The size of the retained data is further refined via pre-training experiments comparing the top 40B, 80B, 120B, and 160B tokens, with the first round prioritizing the top 40B tokens.

Following the initial phase of data collection, a significant portion of mathematical web pages remains unacquired, primarily because the positive examples used to train the fastText model are not sufficiently representative. To address this limitation, we expand the seed corpus by incorporating additional mathematical web sources, thereby enhancing the fastText model's performance. Concretely, we begin by dividing the complete Common Crawl dataset into non-overlapping domains, where each domain comprises web pages sharing a common base URL. For every domain, we determine the fraction of pages retrieved during the first collection cycle. Domains from which more than 10\% of pages were gathered are deemed to be mathematics-oriented (such as \url{mathoverflow.net}).

Next, we manually tag the URLs within these mathematics-focused domains that specifically pertain to mathematical materials (e.g., \url{mathoverflow.net/questions}). Any uncollected pages under these URLs are subsequently added to our seed corpus. This procedure increases the number of positive samples, enabling the fastText model to better recognize mathematical content in subsequent collection cycles. By the end of four collection iterations, we compile a set of 35.5M mathematical web pages, amassing a total of 120B tokens. Upon reaching the fourth round, we observe that roughly 98\% of the pages collected were already included after the third cycle, prompting us to halt further data gathering.

To prevent data leakage from evaluation datasets, we adhere to the filtering protocol outlined in \cite{deepseek-coder}: web pages containing question or answer content from prominent English mathematical benchmarks (such as GSM8K \citep{gsm8k} and MATH \citep{MATH}) and Chinese benchmarks (such as CMATH \citep{wei2023cmath} and AGIEval \citep{agieval}) are systematically removed. The filtering process entails deleting any segment of text that contains an exact 10-gram match with substrings in the evaluation datasets. For cases where benchmark excerpts consist of fewer than 10, but at least 3, grams, we apply exact matching to exclude affected web pages from the mathematical corpus.

\subsection{Validating the Quality of the DeepSeekMath~Corpus}

We conduct pre-training trials to evaluate the DeepSeekMath~Corpus in relation to contemporary math-oriented pre-training datasets:

\begin{itemize}[topsep=0pt]
    \item \textbf{MathPile} \citep{mathpile}: A compilation totaling 8.9B tokens sourced from diverse origins—textbooks, Wikipedia, ProofWiki, CommonCrawl, StackExchange, and arXiv—with arXiv papers constituting the bulk (over 85\%) of the content;
    \item \textbf{OpenWebMath} \citep{openwebmath}: A dataset comprising 13.6B tokens of CommonCrawl data filtered to extract mathematical materials;
    \item \textbf{Proof-Pile-2} \citep{llemma}: A composite dataset combining OpenWebMath, AlgebraicStack (10.3B tokens of mathematical code), and 28.0B tokens of arXiv documents. Experiments utilizing Proof-Pile-2 adopt an arXiv:Web:Code proportion of 2:4:1, following the methodology of \cite{llemma}.
\end{itemize}

\subsubsection{Training Setting}
\label{sec:quality-policy}
We fine-tune a generalist pre-trained language model consisting of 1.3 billion parameters, following the architecture of the DeepSeek LLMs \citep{deepseek-llm}, henceforth referred to as DeepSeek-LLM 1.3B. The model undergoes separate training runs on each mathematical dataset, each comprising 150B tokens. All training procedures leverage the resource-efficient HAI-LLM \citep{haillm} training framework. Consistent with the methodology established in DeepSeek LLMs, the AdamW optimizer \citep{adamW} is employed with $\beta_1=0.9$, $\beta_2=0.95$, and $\mathrm{weight\_decay}=0.1$. A multi-stage learning rate policy is used: the learning rate is linearly ramped up to its peak over the initial 2,000 warmup iterations, subsequently decays to 31.6\% of the peak at 80\% progress through training, and further decreases to 10.0\% of the peak after reaching 90\% completion. The learning rate is capped at 5.3e-4, while a batch size corresponding to 4 million tokens and a context length of 4,000 tokens are maintained throughout training.

\subsubsection{Evaluation Results}

\textbf{The DeepSeekMath~Corpus stands out as a high-quality, multilingual resource and is unrivaled in scale.}

\begin{itemize}[topsep=0pt]
    \item \textbf{High-quality}:  
    We assess downstream capability on eight mathematical reasoning benchmarks utilizing few-shot chain-of-thought prompting \cite{cot}. According to Table \ref{tab:corpora_comparison}, models fine-tuned on the DeepSeekMath~Corpus substantially outperform their counterparts. As depicted in Figure \ref{fig:corpora_comparisons}, the DeepSeekMath-trained model exhibits higher scores than models trained on Proof-Pile-2 when evaluated at 50B tokens (equivalent to one full pass over Proof-Pile-2), signifying a superior average quality of the DeepSeekMath~Corpus.
    \item \textbf{Multilingual}:  
    DeepSeekMath~Corpus incorporates data across several languages, with English and Chinese making up the majority of its content. Results in Table \ref{tab:corpora_comparison} confirm that training on DeepSeekMath~Corpus significantly improves mathematical reasoning in both English and Chinese. By contrast, other widely used corpora, which are largely limited to English, provide little to no benefit—and sometimes a detriment—for Chinese-language mathematical tasks.
    \item \textbf{Large-scale}:  
    The size of the DeepSeekMath~Corpus far exceeds that of prior mathematical datasets. As Figure \ref{fig:corpora_comparisons} illustrates, DeepSeek-LLM 1.3B models trained with DeepSeekMath~Corpus not only exhibit sharper improvements but also sustain progress for longer training periods. In comparison, baseline corpora are considerably smaller, having been seen by the models multiple times during training, which results in quickly saturated performance.
\end{itemize}

\subsection{Training and Evaluating DeepSeekMath-Base 7B}

This section presents DeepSeekMath-Base 7B, a foundational model distinguished by advanced reasoning capabilities, particularly in mathematics. The model initialization utilizes DeepSeek-Coder-Base-v1.5 7B \citep{deepseek-coder}, followed by training on a dataset comprising 500B tokens. The composition of the training data is as follows: DeepSeekMath Corpus constitutes 56\%, AlgebraicStack accounts for 4\%, arXiv makes up 10\%, 20\% is sourced from Github code repositories, and the final 10\% consists of English and Chinese natural language material drawn from Common Crawl. Training largely adheres to the procedure outlined in Section \ref{sec:quality-policy}, with the primary deviations being an increased peak learning rate of 4.2e-4 and an expanded batch size of 10M tokens.

To thoroughly analyze the mathematical proficiency of DeepSeekMath-Base 7B, we evaluate its capacity to autonomously generate comprehensive mathematical solutions, employ external tools for problem-solving, and perform formal theorem verification. Additionally, we examine the broader capabilities of the model, encompassing natural language understanding, logical reasoning, and programming aptitude.

\paragraph{Mathematical Problem Solving with Step-by-Step Reasoning}

To measure the model's effectiveness in mathematical problem solving, DeepSeekMath-Base is assessed using few-shot chain-of-thought prompting \citep{cot} across eight distinct benchmarks in both English and Chinese. These benchmarks feature a range of tasks, including quantitative reasoning (such as GSM8K \citep{gsm8k}, MATH \citep{MATH}, and CMATH \citep{wei2023cmath}) and multiple-choice assessments (for instance, MMLU-STEM \citep{mmlu} and Gaokao-MathQA \citep{agieval}), spanning topics from basic to advanced, covering mathematics at both elementary and collegiate levels.

As detailed in Table \ref{tab:base_math_cot}, DeepSeekMath-Base 7B achieves top results across all eight evaluation benchmarks when compared with open-source base models, encompassing the widely-adopted general-purpose model Mistral 7B \citep{mistral} and the more recent Llemma 34B \citep{llemma}, which received specialized math training via Proof-Pile-2 \citep{llemma}. Of particular note is the substantial margin on the MATH competition-level dataset, where DeepSeekMath-Base outpaces other open-source base models by over 10\% absolute and even surpasses Minerva 540B \citep{minerva}, a much larger proprietary model (77 times the parameter count) based on PaLM \citep{palm} and further refined on mathematical texts.

\paragraph{Mathematical Problem Solving with Tool Use} To assess mathematical reasoning augmented by code execution, we test few-shot program-of-thought prompting \citep{pot,pal} on GSM8K and MATH. The models receive prompts encouraging them to generate Python code for problem-solving, leveraging modules like \textit{math} and \textit{sympy} for advanced calculations. The final program output is treated as the answer. As reported in Table \ref{tab:base_math_pot_proof}, DeepSeekMath-Base 7B sets a new benchmark, outperforming the previous leading model, Llemma 34B.

\paragraph{Formal Mathematics}

Automating formal proof construction improves both the reliability and efficiency of mathematical reasoning, a research focus of growing interest. We test DeepSeekMath-Base 7B in the informal-to-formal proving task outlined by \citep{dsp_proof}, wherein the objective is to construct a formal proof given an informal claim, its formal analogue, and an informal demonstration. Evaluation is conducted on the miniF2F benchmark \citep{minif2f}, which features Olympiad-level formal math problems, prompting the model to generate Isabelle proofs with few-shot examples. Following the approach in \cite{dsp_proof}, the models produce proof outlines, and the Sledgehammer \citep{sledgehammer} prover completes any unresolved steps. Table \ref{tab:base_math_pot_proof} shows that DeepSeekMath-Base 7B achieves notable performance in the domain of proof autoformalization.

\paragraph{Natural Language Understanding, Reasoning, and Code}

We examine the model’s aptitude in natural language understanding using MMLU \citep{mmlu}, its reasoning abilities with BBH \citep{bbh}, and coding performance via HumanEval \citep{codex} and MBPP \citep{mbpp}. Table \ref{tab:understanding_reasoning_code} indicates that DeepSeekMath-Base 7B yields marked gains on MMLU and BBH relative to its predecessor, DeepSeek-Coder-Base-v1.5 \citep{deepseek-coder}, demonstrating the beneficial effects of mathematical training for both language comprehension and logical reasoning. Furthermore, through the inclusion of code-specific tokens during continued training, DeepSeekMath-Base 7B preserves the coding performance previously established by DeepSeek-Coder-Base-v1.5 across both code benchmarks. Across the reasoning and coding benchmarks, DeepSeekMath-Base 7B also significantly exceeds the performance of Mistral 7B \citep{mistral}.

\section{Supervised Fine-Tuning}

\subsection{SFT Data Curation}

We compile an instruction-tuning dataset for mathematics that encompasses both English and Chinese questions, sourced from a variety of mathematical domains and spanning different levels of difficulty. Each question is paired with its corresponding solution, presented in formats such as chain-of-thought (CoT) \citep{cot}, program-of-thought (PoT) \citep{pot,pal}, and tool-integrated reasoning \citep{tora}. The dataset comprises a total of 776K training samples.

\begin{itemize}[topsep=0pt]
    \item \textbf{English mathematical datasets}: Our English-language corpus features GSM8K and MATH, which we annotate with tool-integrated solution strategies, alongside a curated portion of MathInstruct \citep{MathInstruct} and the training partition of Lila-OOD \citep{lila}, with their problems addressed via CoT or PoT. This set includes content from a broad array of mathematical disciplines, such as algebra, probability, number theory, calculus, and geometry.
    \item \textbf{Chinese mathematical datasets}: For Chinese-language problems, we gather K-12 level mathematical questions that encompass 76 sub-categories, including topics like linear equations. Solutions for these are annotated using both CoT and tool-integrated reasoning styles.
\end{itemize}

\subsection{Training and Evaluating DeepSeekMath-Instruct 7B}

This section describes the DeepSeekMath-Instruct 7B model, which is derived from DeepSeekMath-Base and further refined through mathematical instruction tuning. During training, samples are combined in random order up to a cap of 4K tokens for context length. Model training proceeds for 500 steps with a batch size set to 256 and utilizes a fixed learning rate of 5e-5.

To measure mathematical capabilities with and without tool assistance, we assess the model on four quantitative reasoning benchmarks available in both English and Chinese. For reference, we compare its performance with leading contemporary models:

\begin{itemize}[topsep=0pt]
    \item \textbf{Closed-source models} considered are: (1) the GPT series, notably GPT-4 \citep{gpt4} and GPT-4 Code Interpreter~\footnote{\url{https://openai.com/blog/chatgpt-plugins\#\#code-interpreter}}, (2) Gemini Ultra and Pro \citep{gemini}, (3) Inflection-2 \citep{inflection-2}, (4) Grok-1~\footnote{\url{https://x.ai/model-card}}, as well as (5) Baichuan-3~\footnote{\url{https://www.baichuan-ai.com}}, and (6) the most recent GLM-4~\footnote{\url{https://open.bigmodel.cn/dev/api\#glm-4}} from the GLM family \citep{glm}, which are the latest models launched by major Chinese firms.
\end{itemize}

These models are designed for broad applicability, and the majority have been subjected to alignment processes.  
\item \textbf{Open-source models} comprise both general-purpose systems such as (1) DeepSeek-LLM-Chat 67B \citep{deepseek-llm}, (2) Qwen 72B \citep{qwen}, (3) SeaLLM-v2 7B \citep{seallm}, and (4) ChatGLM3 6B \citep{chatglm3}, alongside math-specialized models, including: (5) InternLM2-Math 20B~\footnote{\url{https://github.com/InternLM/InternLM-Math}}, which extends InternLM2 through dedicated math training followed by instruction fine-tuning; (6) Math-Shepherd-Mistral 7B, which incorporates PPO optimization \citep{schulman2017proximal} into Mistral 7B \citep{mistral} using a reward model guided by process supervision; (7) the WizardMath series \citep{wizardmath}, enhancing mathematical abilities of Mistral 7B and Llama-2 70B \citep{llama2} via evolve-instruct (an instruction tuning variant leveraging instructions evolved by AI) and PPO, with data primarily from GSM8K and MATH; (8) MetaMath 70B \citep{metamath}, fine-tuned from Llama-2 70B on an augmented dataset that incorporates GSM8K and MATH; (9) ToRA 34B \cite{tora}, which adapts CodeLlama 34B for mathematical reasoning integrated with external tools; and (10) MAmmoTH 70B \citep{MathInstruct}, which applies instruction tuning with MathInstruct to Llama-2 70B.
\end{itemize}

According to Table \ref{tab:sft_rl_math}, when models are restricted from utilizing external tools, DeepSeekMath-Instruct 7B delivers outstanding performance in sequential reasoning. Remarkably, on the advanced MATH benchmark, this model exceeds all other open-source alternatives and outperforms most proprietary models—including Inflection-2 and Gemini Pro—by no less than a 9\% absolute margin. This holds true even against models of significantly greater scale (e.g., Qwen 72B) or those trained with math-centric reinforcement learning enhancements (e.g., WizardMath-v1.1 7B). Although DeepSeekMath-Instruct achieves parity with Chinese proprietary models such as GLM-4 and Baichuan-3 on MATH, it continues to lag behind GPT-4 and Gemini Ultra.

In evaluation scenarios allowing integration of natural language and program-based tools, DeepSeekMath-Instruct 7B attains nearly 60\% accuracy on MATH, outperforming all available open-source competitors. Across additional benchmarks, its performance is comparable to DeepSeek-LLM-Chat 67B, which was previously the best-performing open-source model and is ten times larger.

\section{Reinforcement Learning}

\subsection{Group Relative Policy Optimization} Reinforcement learning (RL) has consistently demonstrated the capacity to further strengthen mathematical reasoning abilities in LLMs beyond the Supervised Fine-Tuning (SFT) phase \citep{wang2023math,wizardmath}. Here, we present Group Relative Policy Optimization (GRPO), our proposed RL method designed to be both efficient and effective.

\subsubsection{From PPO to GRPO} Proximal Policy Optimization (PPO) \citep{schulman2017proximal} is a prominent actor-critic RL approach that serves as a cornerstone for the RL fine-tuning phase in LLM development \citep{ouyang2022training}. PPO seeks to refine LLMs by optimizing the following surrogate objective:

\begin{equation}
\footnotesize
    \mathcal{J}_{PPO}(\theta) = \mathbb{E}{[q \sim P(Q), o \sim \pi_{\theta_{old}}(O|q)]} \frac{1}{|o|} \sum_{t=1}^{|o|} \min \left[ \frac{\pi_\theta(o_{t} | q, o_{<t})}{\pi_{\theta_{old}}(o_{t} | q, o_{<t})} A_{t}, \text{clip} \left( \frac{\pi_\theta(o_{t} | q, o_{<t})}{\pi_{\theta_{old}}(o_{t} | q, o_{<t})}, 1 - \varepsilon, 1 + \varepsilon \right)  A_{t

Here, $\pi_{\theta}$ denotes the current policy model, while $\pi_{\theta_{old}}$ refers to the preceding policy. The variables $q$ and $o$ correspond to questions and their associated outputs, which are sampled from a dataset of questions and generated using $\pi_{\theta_{old}}$, respectively. The hyper-parameter $\varepsilon$ plays a role in the clipping mechanism introduced by PPO to ensure stable learning. The advantage $A_t$ is estimated through Generalized Advantage Estimation (GAE) \citep{gae}, leveraging reward sequences $\{r_{\ge t}\}$ and predictions from a learned value function $V_{\psi}$. Accordingly, in PPO, simultaneous training of both a value function and a policy model is required. To reduce reward model over-optimization, it is standard practice to incorporate a per-token Kullback-Leibler (KL) divergence penalty—based on a reference policy—directly into the reward for each token \citep{ouyang2022training}, specifically as follows:

\begin{equation}
    r_{t} = r_\varphi(q, o_{\le t}) - \beta \log\frac{\pi_{\theta}(o_{t}|q, o_{<t})}{\pi_{ref}(o_{t}|q, o_{<t})},
\label{eq:PPO-reward}
\end{equation}

Here, $r_\varphi$ denotes the reward function, $\pi_{ref}$ represents a reference policy—frequently taken as the starting SFT model—and $\beta$ sets the strength of the KL penalty term.

Since PPO employs a value function that is often architecturally as complex as the policy model itself, this results in considerable computational and memory overhead. Furthermore, the value function's role during reinforcement learning is as a variance-reducing baseline in the computation of advantages. However, in large language model scenarios, reward models typically only supply a reward for the final token, which complicates accurate value estimation at each token. To overcome these issues, and as depicted in Figure \ref{fig:grpo}, we introduce Group Relative Policy Optimization (GRPO). GRPO dispenses with an explicit value function, unlike PPO, by instead utilizing the mean reward across multiple output samples—each in response to the same question—as the baseline. Specifically, for each prompt $q$, GRPO generates a set of outputs $\{o_1, o_2, \cdots, o_G\}$ from $\pi_{\theta_{old}}$, and then trains the policy by maximizing the objective:

\begin{equation}
\footnotesize
\begin{split}
    \mathcal{J}_{GRPO}(\theta) &= \mathbb{E}{[q \sim P(Q), \{o_i\}_{i=1}^G \sim \pi_{\theta_{old}}(O|q)]}  \\
    & \frac{1}{G}\sum_{i=1}^G\frac{1}{|o_i|} \sum_{t=1}^{|o_i|} \left\{ \min \left[ \frac{\pi_\theta(o_{i,t} | q, o_{i,<t})}{\pi_{\theta_{old}}(o_{i,t} | q, o_{i,<t})} \hat{A}_{i,t}, \text{clip} \left( \frac{\pi_\theta(o_{i,t} | q, o_{i,<t})}{\pi_{\theta_{old}}(o_{i,t} | q, o_{i,<t})}, 1 - \varepsilon, 1 + \varepsilon \right)  \hat{A}_{i,t} \right] - \beta \mathbb{D}_{KL}\left[\pi_{\theta} || \pi_{ref}\right]\right\} ,
\end{split}
\label{eq:GRPO-obj}
\end{equation}

In this expression, $\varepsilon$ and $\beta$ continue to act as hyper-parameters, and $\hat{A}_{i,t}$ signifies the advantage calculated using relative rewards among outputs within the same group—a procedure detailed in subsequent subsections. This group-relative computation of advantages resonates with the design of many reward models, which are typically trained using comparison datasets consisting of multiple outputs for a shared question. Additionally, instead of incorporating the KL penalty within the reward itself, GRPO imposes regularization by directly appending the KL divergence between the candidate and reference policies to the objective function, thus simplifying the process of determining $\hat{A}_{i,t}$. Distinct from the KL penalty formulation in (\ref{eq:PPO-reward}), our approach estimates KL divergence using the following unbiased estimator \cite

\begin{algorithm}[t]
  \small
  \caption{Iterative Group Relative Policy Optimization}
  \textbf{Input} initial policy model $\pi_{\theta_{\text{init}}}$; reward models $r_\varphi$; task prompts $\mathcal{D}$; 
  hyperparameters $\varepsilon$, $\beta$, $\mu$
  \begin{algorithmic}[1]
    \State Assign initial policy model $\pi_\theta \leftarrow \pi_{\theta_{\text{init}}}$
    \For{iteration = 1, \dots, I}
      \State Set reference model $\pi_{ref} \leftarrow \pi_{\theta}$
      \For{step = 1, \dots, M}
        \State Draw a batch $\mathcal{D}_b$ from $\mathcal{D}$
        \State Update the previous policy model $\pi_{\theta_{old}} \leftarrow \pi_{\theta}$ 
        \State For each question $q \in \mathcal{D}_b$, generate $G$ outputs $\{o_i\}_{i=1}^G$ by sampling from $\pi_{\theta_{old}} (\cdot \mid q) $
        \State Evaluate sampled outputs $\{o_i\}_{i=1}^G$ with $r_\varphi$ to obtain rewards $\{r_i\}_{i=1}^{G}$
        \State Calculate $\hat{A}_{i,t}$ for the $t$-th token in $o_i$ using group relative advantage estimation
        \For{GRPO iteration = 1, \dots, $\mu$}
          \State Update policy $\pi_{\theta}$ by optimizing the GRPO objective (Equation \ref{eq:GC-GRPO})
        \EndFor
      \EndFor 
      \State Continue training $r_\varphi$ using a replay strategy.
    \EndFor 
  \end{algorithmic}
  \textbf{Output} $\pi_\theta$
  \label{alg:iter-grpo}
\end{algorithm}

\subsubsection{Outcome Supervision RL with GRPO} Specifically, for a given question $q$, the old policy $\pi_{\theta_{old}}$ generates a set of candidate outputs $\{o_1, o_2, \cdots, o_G\}$. Each output is evaluated by the reward model, resulting in a collection of rewards $\mathbf{r}=\{r_1, r_2, \cdots, r_G\}$. The rewards are then standardized by centering with the group mean and scaling by the group standard deviation. In outcome supervision, this standardized value is assigned as the terminal reward to each output $o_i$ and used as the advantage $\hat{A}_{i, t}$ for all tokens within that output; in other words, $\hat{A}_{i, t} = \widetilde{r}_i = \frac{r_i- {\rm mean}(\mathbf{r})}{{\rm std}(\mathbf{r})}$. Policy optimization proceeds by maximizing the objective outlined in equation (\ref{eq:GRPO-obj}).

\subsubsection{Process Supervision RL with GRPO} While outcome supervision offers a reward only at the completion of each output, this can be limiting for supervising complex multi-step reasoning problems. In line with \cite{wang2023math}, we also employ process supervision, whereby a reward is provided following every reasoning step. Formally, with $q$ and generated outputs $\{o_1, o_2, \cdots, o_G\}$, a process-based reward model evaluates each intermediate step, resulting in reward sets: $\mathbf{R} = \{ \{r_1^{index(1)},\cdots,r_1^{index(K_1)}\}, \cdots,  \{r_G^{index(1)},\cdots,r_G^{index(K_G)}\} \}$, where $index(j)$ refers to the end token of step $j$ and $K_i$ denotes the total steps in output $i$. All such stepwise rewards are normalized as $\widetilde{r}_i^{index(j)} = \frac{r_i^{index(j)} - {\rm mean(\mathbf{R})}}{{\rm std(\mathbf{R})}}$. Then, the advantage $\hat{A}_{i, t}$ assigned to the $t$-th token is calculated by summing all normalized future rewards across subsequent steps: $\hat{A}_{i, t} = \sum_{index(j) \ge t} \widetilde{r}_i^{index(j)}$. Finally, optimization is performed to maximize the policy objective presented

\subsubsection{Iterative RL with GRPO} During the progression of reinforcement learning, the previously trained reward model may no longer provide adequate supervision for the evolving policy model. Consequently, we investigate the application of iterative RL with GRPO. As depicted in Algorithm \ref{alg:iter-grpo}, the iterative GRPO procedure involves regenerating datasets for the reward model, drawing samples from the current policy model. The reward model is continuously updated through a replay mechanism that retains 10\% of prior data. Subsequently, the reference model is synchronized with the policy model, and the policy model undergoes continual training leveraging the updated reward model.

\subsection{Training and Evaluating DeepSeekMath-RL}

We implement reinforcement learning using DeepSeekMath-Instruct 7B. The RL training data comprises chain-of-thought-style questions derived from GSM8K and MATH subsets in the SFT dataset, amounting to approximately 144K instances. To specifically assess the influence of RL on benchmarks where relevant data is not provided during RL, other SFT data are omitted. The training data for the reward models are constructed according to the protocol in \citep{wang2023math}. We initialize the reward model on DeepSeekMath-Base 7B, setting the learning rate to 2e-5. For the GRPO approach, the policy model utilizes a learning rate of 1e-6, and the KL coefficient is set at 0.04. Each prompt is associated with $64$ sampled completions. The sequence length limit is 1024 tokens, and the batch size for training is also 1024. Each exploration stage is followed by a single policy model update. DeepSeekMath-RL 7B is assessed on the same benchmarks as DeepSeekMath-Instruct 7B. In the context of DeepSeekMath-RL 7B, tasks from GSM8K and MATH with chain-of-thought processes are categorized as in-domain, while all remaining benchmarks are treated as out-of-domain.

Table \ref{tab:sft_rl_math} presents comparative results for both open- and closed-source models engaging in chain-of-thought and tool-integrated reasoning across English and Chinese datasets. Our analysis reveals the following insights: (1) Utilizing chain-of-thought reasoning, DeepSeekMath-RL 7B achieves accuracies of 88.2\% on GSM8K and 51.7\% on MATH. These results exceed those of any other open-source models within the 7B–70B parameter spectrum and outperform most closed-source alternatives as well. (2) Notably, DeepSeekMath-RL 7B is trained exclusively on chain-of-thought-formatted instruction tuning data from GSM8K and MATH, with DeepSeekMath-Instruct 7B as its initialization. Although its training data is comparatively limited, DeepSeekMath-RL 7B consistently surpasses DeepSeekMath-Instruct 7B across all assessment criteria, thereby highlighting the advantages conferred by reinforcement learning.

\section{Discussion}
This section details our observations derived from pre-training and reinforcement learning experiments.

\subsection{Insights from Pre-Training}

We begin by sharing the lessons gleaned from pre-training experiments. Unless indicated otherwise, experimental conditions correspond to those in Section \ref{sec:quality-policy}. The DeepSeekMath Corpus discussed herein refers specifically to an 89-billion-token version collected in the data gathering process's second phase.

\subsubsection{Code Training as a Catalyst for Mathematical Reasoning}

Although the notion that code training enhances reasoning is widely held, empirical evidence remains sparse. Our experiments provide partial support for this idea in mathematical contexts: incorporating code during training elevates the model's mathematical reasoning, independent of whether external tools are used.

To evaluate the effect of code training on mathematical reasoning capabilities, we contrasted two different training regimes: a sequential (two-stage) approach and a unified (one-stage) training scheme.

\textbf{Two-Stage Training}

\begin{itemize}[topsep=0pt]
    \item \textbf{Code Training for 400B Tokens $\rightarrow$ Math Training for 150B Tokens}: We initially pre-train DeepSeek-LLM 1.3B on 400B code tokens and then further fine-tune the model on an additional 150B math tokens;
    \item \textbf{General Training for 400B Tokens $\rightarrow$ Math Training for 150B Tokens}: To serve as a baseline, we conduct an alternative experiment in which DeepSeek-LLM 1.3B undergoes pre-training with 400B general tokens (sampled from DeepSeek-AI’s extensive general-purpose corpus) in place of code tokens, subsequently followed by math token fine-tuning, so as to compare the effects of initial code versus general data on downstream mathematical abilities.
\end{itemize}

\textbf{One-Stage Training}

\begin{itemize}[topsep=0pt]
    \item \textbf{Math Training for 150B Tokens}: Here, DeepSeek-LLM 1.3B is trained solely with 150B tokens from the math domain;
    \item \textbf{Training on a mixture of 400B Code Tokens and 150B Math Tokens}: Noting that sequential math training after code pre-training can negatively impact previously acquired coding skills, we also explore a single-stage approach in which math and code tokens (totaling 550B) are mixed throughout training to test if this simultaneous exposure bolsters mathematical performance and preserves coding abilities, thereby reducing catastrophic forgetting.
\end{itemize}

\paragraph{Results}
The outcomes across the evaluated configurations are presented in Table \ref{tab:code-math} and Table \ref{tab:code-math-general-eval}.

In both staged and mixed training paradigms, incorporating code token pre-training enhances program-supported mathematical reasoning. From Table \ref{tab:code-math}, code-first training in a two-stage protocol leads to substantial gains in solving GSM8K and MATH benchmarks when models are permitted Python usage, with further refinements obtained by subsequent math token training. Notably, the one-stage joint code/math training scenario addresses the challenge of catastrophic forgetting that arises in sequential (two-stage) setups, while simultaneously promoting both coding performance (Table \ref{tab:code-math-general-eval}) and code-assisted mathematical problem-solving (Table \ref{tab:code-math}).

Moreover, code pre-training also boosts mathematical reasoning in non-tool (purely mathematical) settings. Sequential code-math training yields incremental improvements from the initial code stage and accelerates the learning effect during math fine-tuning, cumulatively achieving the highest scores observed. Conversely, joint code/math training within a single-stage regime results in diminished mathematical reasoning (in non-tool contexts). This limitation likely stems from the constrained model scale—DeepSeek-LLM 1.3B may not possess sufficient capacity to absorb and integrate both modalities in a concurrent training schedule.

\subsubsection{ArXiv Papers Appear to Have Limited Effect on Mathematical Reasoning}

In the field, arXiv papers are routinely incorporated into the pre-training datasets for mathematical large language models \citep{minerva,gpt-f,llemma,mathpile}. Nevertheless, there has been little comprehensive assessment of their true influence on models' mathematical reasoning capacities. Surprisingly, our experimental findings indicate that incorporating arXiv papers does not lead to substantial enhancements in mathematical reasoning abilities. To evaluate this, we trained models with different parameter scales, specifically DeepSeek-LLM 1.3B and DeepSeek-Coder-Base-v1.5 7B \citep{deepseek-coder}, utilizing distinct arXiv-derived datasets with diverse preprocessing workflows:

\begin{itemize}[topsep=0pt]
    \item \textbf{MathPile} \citep{mathpile}: consisting of 8.9B tokens gathered with both cleaning and heuristic-based filtering methods, where scientific arXiv articles make up more than 85\% of the content;
    \item \textbf{ArXiv-RedPajama} \citep{redpajama}: comprising 28.0B tokens extracted from all arXiv LaTeX files, which have had extraneous elements such as preambles, comments, macros, and bibliographies excluded.
\end{itemize}

For this investigation, DeepSeek-LLM 1.3B was trained for 150B tokens, while DeepSeek-Coder-Base-v1.5 7B was trained for 40B tokens on each of the arXiv datasets independently. The results reveal that exposure solely to arXiv data does not foster better mathematical reasoning. Rather, models trained exclusively on such corpora show either stagnation or decline on a broad array of mathematical evaluation suites of varied difficulty used in this work. These assessments range from quantitative reasoning datasets (such as GSM8K and MATH, shown in Table \ref{tab:arxiv-cot}), to multiple-choice tasks like MMLU-STEM (see Table \ref{tab:arxiv-cot}), and formal mathematics benchmarks, notably miniF2F (Table \ref{tab:arxiv_atp}).

It is important to acknowledge the limitations of this analysis, which mean the findings should not be considered definitive. Specifically, this work has not yet addressed:

\begin{itemize}[topsep=0pt]
    \item Potential influence of arXiv data on specialized mathematical tasks omitted from this research, for instance, translating formal mathematics into more informal descriptions;
    \item Effects arising when arXiv data is mixed with additional data modalities or sources;
    \item The possibility that arXiv data may yield improvements at much larger model scales than those tested here.
\end{itemize}

Accordingly, these questions remain open and merit further examination in future research efforts.

\subsection{Insights on Reinforcement Learning}
\subsubsection{A Move Toward a Unified Training Framework} In this section, we introduce a unifying framework to compare a spectrum of training algorithms, encompassing SFT, RFT, DPO, PPO, GRPO, and proceed to empirically investigate the parameters and components underpinning this framework. In general terms, the gradient with respect to the parameters $\theta$ of a training algorithm is given by:

\begin{equation}
    \nabla_{\theta}\mathcal{J}_{\textcolor{red}{\mathcal{A}}}(\theta) = \mathbb{E}[\underbrace{(q,o) \sim \textcolor{red}{\mathcal{D}}}_{Data \ Source}]\left( \frac{1}{|o|} \sum_{t=1}^{|o|}  \underbrace{GC_{{\mathcal{A}}}(q, o, t, \textcolor{red}{\pi_{{rf}}})}_{Gradient \ Coefficient}  \nabla_{\theta}\log \pi_{\theta}(o_t | q, o_{<t})\right).
\label{eq:objective}
\end{equation}

This formulation encompasses three primary components: 1) the \textit{Data Source} $\mathcal{D}$, which specifies the collection of samples for training; 2) the \textit{Reward Function} $\pi_{{rf}}$, which supplies the supervisory reward signals; and 3) the \textit{Algorithm} $\mathcal{A}$, which synthesizes information from both the data and the reward to yield the gradient coefficient $GC$—thereby calibrating the scale of positive or negative updates assigned to each data point. We examine a range of influential techniques within this overarching framework:

\begin{itemize}
    \item  \textbf{Supervised Fine-tuning (SFT)}: SFT involves additional training of a pretrained model using datasets curated by humans.
    \item \textbf{Rejection Sampling Fine-tuning (RFT)}: In RFT, the SFT model undergoes further training with outputs, filtered for correctness, that are generated in response to the original SFT prompts.
    \item \textbf{Direct Preference Optimization (DPO)}: DPO enhances the SFT model by utilizing pairwise losses and supplementary samples, leading to improved output alignment with preferences.
    \item \textbf{Online Rejection Sampling Fine-tuning (Online RFT)}: Diverging from standard RFT, Online RFT continually updates the policy model by drawing and refining training outputs generated on-the-fly by the current policy.
    \item \textbf{PPO/GRPO}: Starting from an SFT-initialized policy, PPO/GRPO further develops the model using online sampling of outputs, with training guided by real-time feedback from the evolving policy.
\end{itemize}

The elements constituting these approaches are outlined in Table \ref{tab:data_policy}. For an in-depth derivation, readers are directed to Appendix \ref{app:analysis-rl}.

\paragraph{Observation about Data Source} We categorize the data sources into two main types: online and offline sampling. Online sampling refers to training datasets generated dynamically from the policy model's live exploration, whereas offline sampling utilizes data gathered from the initial SFT model. Methods such as RFT and DPO adopt the offline approach, while Online RFT and GRPO rely on online data collection.

In Figure \ref{fig:rl-analysis}, it is evident that Online RFT surpasses RFT across both evaluated benchmarks. During the early training phase, Online RFT performs on par with RFT; however, as training progresses, Online RFT exhibits a clear superiority, underlining the benefits of online data acquisition. This outcome is expected: in the initial epochs, the actor model's behavior closely mirrors that of the SFT model, resulting in sampled data with only marginal variations. As training advances, the actor's generated samples increasingly diverge, and the advantage of continuously updating the training data with real-time samples becomes more pronounced.

\paragraph{Observation about Gradient Coefficient} The conversion of the reward signal into a gradient coefficient for model updates varies by algorithm. In our work, we use two forms of reward functions: `Rule' and `Model'. The 'Rule' metric is based on whether a response's answer is correct, while 'Model' involves training a reward model to assign scores to responses, with the reward model’s data annotated according to rule-based judgments. As indicated by Equations \ref{eq:GC-RFT} and \ref{eq:GC-GRPO}, GRPO distinguishes itself from Online RFT by modulating its gradient coefficient in accordance with the specific reward values predicted by the reward model. This design permits more granular positive and negative reinforcement aligned with the magnitude of the rewards. On the other hand, Online RFT employs a uniform reinforcement strategy: incorrect responses are not penalized, and all correct answers are reinforced equally, regardless of their degree of correctness.

As illustrated in Figure \ref{fig:rl-analysis}, GRPO consistently outperforms online RFT, underscoring the advantage of modulating positive and negative gradient coefficients. Moreover, GRPO+PS achieves better results than GRPO+OS, underscoring the gains afforded by more granular, step-wise gradient coefficient adjustments. Additionally, our study extends to iterative RL; two iterative rounds were conducted in our experiments. The results, depicted in Figure \ref{fig:iter-rl}, reveal marked performance gains, particularly after the initial iteration.

\subsubsection{Why RL Works?}  
Within this study, we employ reinforcement learning leveraging a subset of instruction tuning datasets, observing marked improvements relative to the instruction-tuned baseline. To elucidate the mechanisms driving RL’s efficacy, we analyze Pass@K and Maj@K accuracies for both Instruct and RL models across two evaluation benchmarks. According to Figure \ref{fig:maj-pass}, reinforcement learning raises Maj@K but not Pass@K. This suggests RL strengthens model robustness mainly by increasing the likelihood that the correct response appears among the top-K outputs, rather than by advancing intrinsic model capabilities. \textbf{The data suggest performance gains are due to amplifying the rank of the correct output within TopK responses rather than substantive improvements in the core abilities.} In a related vein, \citep{wang2023making} discussed a \textbf{misalignment problem} in the reasoning abilities of SFT models and demonstrated that preference alignment strategies can notably enhance their reasoning outcomes \citep{yuan2023rrhf,song2023preference,wang2023making}.

\subsubsection{How to Achieve More Effective RL?}
\label{sec:effec_rl}
Our findings show that reinforcement learning achieves strong results on mathematical reasoning benchmarks. We further propose a unified analytical framework that brings together various well-known training methodologies under the common umbrella of either direct or simplified RL approaches. As depicted in Equation \ref{eq:objective}, three fundamental elements emerge: Data Source, Algorithm, and Reward Function. Below, we outline possible future directions regarding each.

\paragraph{Data Source} The data source serves as the fundamental input for all training paradigms. In RL contexts, this specifically refers to unlabeled problems along with model-generated candidate solutions sampled from the policy. Here, we limited ourselves to question sets used during instruction tuning, employing straightforward nucleus sampling for response generation. We hypothesize that this constraint partially explains why RL mainly advances Maj@K performance in our pipeline. Moving forward, we intend to expand the RL setup to encompass out-of-distribution queries and to investigate the utility of \textbf{advanced sampling (decoding) strategies}, such as tree-search-based algorithms \citep{yao2023tree}. Additionally, \textbf{efficient inference techniques} \citep{xia-etal-2023-speculative,leviathan2023fast,kwon2023efficient,xia2024unlocking}, which influence the efficiency of exploration within policy models, are expected to play a vital role in subsequent developments.

\paragraph{Algorithms} Algorithms utilize data and the reward signal to generate gradient coefficients, which are then employed to adjust the model parameters. Referring to Equation \ref{eq:objective}, virtually all approaches currently place considerable \textbf{TRUST} in the reward function’s signal to modify the conditional likelihood of specific tokens, either amplifying or reducing it as needed. Nevertheless, one cannot guarantee that the reward signal is consistently dependable, particularly when facing highly intricate tasks. For instance, the PRM800K datasets \citep{lightman2023let}, even though meticulously labeled by skilled annotators, still exhibit a significant rate of error, with about 20\% of annotations being inaccurate\footnote{\url{https://github.com/openai/prm800k/issues/12\#issuecomment-1728491852}}. Consequently, we aim to investigate reinforcement learning algorithms that can maintain robustness under conditions where reward signals are noisy. We argue that alignment approaches inspired by the \textbf{WEAK-TO-STRONG} paradigm \citep{burns2023weak} have the potential to drive transformative shifts in learning algorithm design.

\paragraph{Reward Function} The reward function constitutes the fundamental source of supervision for training. Within reinforcement learning, this function is typically instantiated as a neural reward model. We identify three key avenues for furthering reward model research: 1) \textbf{Improving the reward model’s generalization capabilities.} Reward models must generalize effectively to tackle novel questions and unconventional decoding behaviors; otherwise, reinforcement learning may simply preserve the current distribution of LLMs, without leading to meaningful advances in their underlying competencies; 2) \textbf{Capturing and representing uncertainty in the reward model.} Proper quantification of uncertainty could facilitate interactions between weak reward models and weak-to-strong alignment algorithms; 3) \textbf{Developing high-quality process reward models efficiently,} enabling the delivery of granular training signals that target specific steps in complex reasoning tasks \citep{lightman2023let,wang2023math}.

\section{Conclusion, Limitation, and Future Work} This paper introduces DeepSeekMath, which establishes a new state-of-the-art among open-source models on the MATH benchmark for competition-level mathematics, reaching performance comparable to proprietary systems. DeepSeekMath~builds on DeepSeek-Coder-v1.5 7B as its foundation, followed by sustained training over 500B tokens, notably including 120B mathematics-specific tokens extracted from Common Crawl. Our comprehensive ablation analysis demonstrates that web-derived content represents a highly valuable resource for curating mathematical datasets, whereas content from arXiv provides less benefit than anticipated. Additionally, we present Group Relative Policy Optimization (GRPO), an adaptation of Proximal Policy Optimization (PPO), capable of enhancing mathematical reasoning abilities while being more memory-efficient. Experimental findings suggest that GRPO yields clear benefits, even for DeepSeekMath-Instruct 7B models already achieving strong benchmark scores. Further, we propose a cohesive framework that synthesizes multiple reinforcement learning techniques, outlining promising directions for advancing their effectiveness.

While DeepSeekMath~achieves strong results on tasks involving quantitative reasoning, its performance in geometric reasoning and theorem-proving remains below that of closed models. For example, our dry-run evaluations reveal the model struggles with triangle and ellipse-related problems, which could reflect biases introduced during data curation and fine-tuning stages. Moreover, the comparatively smaller model size results in inferior few-shot learning capabilities relative to GPT-4: GPT-4 benefits considerably from few-shot prompts, whereas DeepSeekMath~exhibits minimal difference between zero-shot and few-shot performance. Looking ahead, our priorities include further refining our data engineering pipeline to assemble higher-quality pre-training corpora and pursuing the prospective research avenues identified in Section \ref{sec:effec_rl} for more robust reinforcement learning strategies for large language models.

\end{document}